% Full-paper adaptation; the separately reviewed V-C fragment remains unchanged.
% Requires amsmath, amssymb, graphicx, booktabs, algorithm, algpseudocode.
\providecommand{\vcassetroot}{../evidence/05-adam/c-intervention}
\subsection{Interventions Guided by Functional Motion}
\label{sec:vc-interventions}

The tracing analysis separates the direction and magnitude of representation
changes. We therefore test whether limiting local displacement can reduce loss
excursions while preserving training progress and allowing reversals. We compare
a local displacement cap with immediate LR reduction, late second-moment
changes, and global gradient clipping.

\begin{figure*}[!t]
\centering
\includegraphics[width=\textwidth]{\vcassetroot/figures/24-main-intervention-tradeoff.pdf}
\caption{Loss variation and training progress for all 33 treated continuations,
each compared with its unchanged Adam control at the origin. Horizontal values
are $100(\bar L/\bar L_{\mathrm{Adam}}-1)$; vertical values are
$100(1-\mathrm{PV}/\mathrm{PV}_{\mathrm{Adam}})$. Leftward and upward changes are
favorable. Axes are model-specific and retain every setting, including increased
variation. This is one initialization per model; settings have unequal effective
strength and are not independent seed replications.}
\label{fig:vc-tradeoff}
\end{figure*}

\input{\vcassetroot/tables/matched-main-outcomes.tex}

\begin{figure*}[!t]
\centering
\includegraphics[width=\textwidth]{\vcassetroot/figures/25-main-q50-functional-motion.pdf}
\caption{Full training loss and fixed-probe motion under unchanged Adam and the
preselected q50 cap. Heatmaps show consecutive-change cosine and displacement
RMS at five MLP and six ResNet/Transformer stage endpoints. Each incoming change
is placed at its completed-step endpoint. Both RMS maps divide by the same
stage-specific maximum under unchanged Adam; the color scale includes every
value. Dashed vertical lines mark the baseline's largest relative loss-excursion
peak within the full window, not treatment triggers. Embed., Attn and Pool denote
embedding, attention output and pooling endpoints. Output motion uses raw MLP
coordinates or cross-entropy log probabilities. No smoothing or stage averaging
is applied.}
\label{fig:vc-motion}
\end{figure*}

\subsubsection{A past-calibrated displacement cap}
For an affine module $Z=HW^\top+\mathbf{1}_n b^\top$, let
$H_s\in\mathbb{R}^{n\times d}$ contain its current fixed-probe inputs,
$W_s\in\mathbb{R}^{m\times d}$ its weights, and $b_s\in\mathbb{R}^m$ its biases.
Rows index examples or token/spatial positions; columns of $Z\in\R^{n\times m}$
index output units. Let $\Delta W_s$ and $\Delta b_s$ be the native stored Adam increments
with the shapes of $W_s$ and $b_s$. Convolution inputs are receptive-field
patches; absent biases have zero increment.
The direct change attributable to this module's update is
\begin{equation}
 \Delta Z_s^{\mathrm{dir}}=H_s\Delta W_s^\top+\mathbf{1}_n \Delta b_s^\top,\qquad
 c_{s,j}=\frac{\|\Delta Z^{\mathrm{dir}}_{s,:,j}\|_2}{\sqrt{n}}.
 \label{eq:vc-score}
\end{equation}
Thus $c_{s,j}$ is one RMS score per output unit across the probe rows,
rather than a worst-case bound for each example. For a positive module budget
$\tau_\ell$, we scale each complete output-unit
weight row and its bias together:
\begin{equation}
 \begin{split}
 a_{s,j}&=\min(1,\tau_\ell/c_{s,j}),\\
 \Delta W^{\mathrm{cap}}_{s,j:}&=a_{s,j}\Delta W_{s,j:},\qquad
 \Delta b^{\mathrm{cap}}_{s,j}=a_{s,j}\Delta b_{s,j}.
 \end{split}
 \label{eq:vc-cap}
\end{equation}
A zero score uses factor one; a zero budget disables that module's cap.
Before treatment at step $f\to f+1$, one scalar budget per module is fixed as
\begin{equation}
 \begin{split}
 \tau_\ell(q)&=\operatorname{Quantile}_q\{c^{\mathrm{native},\ell}_{s,j}\},\\
 &s=f-8,\ldots,f-1,\quad j=1,\ldots,m_\ell.
 \end{split}
 \label{eq:vc-calibration}
\end{equation}
Quantiles use linear interpolation of the pooled scores. The median (q50)
is preselected; q25/q75 test sensitivity. Units whose proposal would induce
larger direct movement are damped proportionally, while smaller proposals
are retained. The budget has activation units and remains fixed as the current
inputs change.

The rule bounds direct affine RMS on the current probe, up to FP32 storage
rounding. Incoming upstream motion and its interaction with the update remain
unconstrained, so total representation motion and loss are not bounded.
Algorithm~\ref{alg:vc-cap} retains native Adam state and untargeted increments.
Scores use a current probe forward to capture $H_s$ and local products with
the proposed increments; each module does not require a second full-network
forward at altered weights. Evaluating propagated motion separately compares
successive probe captures. The controller requires neither temporal SVD nor
a matrix inverse, but its extra forward, products, and bookkeeping have measured
cost. Supplementary Section S5 gives the input-metric interpretation and the
unconstrained upstream terms.

\begin{algorithm}[!b]
\caption{Past-calibrated local displacement cap}
\label{alg:vc-cap}
\begin{algorithmic}[1]
\Require State $\theta_s$, Adam state, fixed probe, targets, fixed budgets $\tau_\ell$
\State Compute the full objective gradient at $\theta_s$.
\State Capture current probe inputs $H_s^\ell$ with the fixed normalization context; preserve random state.
\State Advance native Adam moments/clock and retain the actual stored proposal increments.
\For{each target $\ell$ with $\tau_\ell>0$}
  \State Compute scores by~\eqref{eq:vc-score}; scale weight rows and biases by~\eqref{eq:vc-cap}.
\EndFor
\State Store $\theta_{s+1}$ with untargeted increments unchanged; retain the advanced Adam state.
\State Check direct-score bounds with storage-rounding tolerance; measure propagated motion and loss separately.
\end{algorithmic}
\end{algorithm}

\subsubsection{Matched comparison}
Within each model, all settings share complete model/optimizer/random state,
data order and numerical settings at one initialization. Fixed windows are
$512$--$800$ for CIFAR-10/MLP,
$256$--$576$ for CIFAR-10/ResNet-32, and $544$--$928$ for the two-layer
WikiText-2 Transformer. Training step $s$ counts completed optimizer updates.
Each update uses the complete objective in Table~\ref{tab:setup}, with its
original normalization groups and no dropout.

Unchanged Adam~\cite{kingma2014adam} uses learning rates $.001/.003/.003$,
$\beta_1=.9$, $\beta_2=.999$, $\epsilon=10^{-8}$, and no weight decay.
The four intervention families are immediate LR multipliers $.75/.5/.1$;
late $\beta_2$ values $.95/.9/.5$; global raw-gradient L2 clipping at
$.5/1$~\cite{pascanu2013difficulty}; and distributed q25/q50/q75 caps.
The $\beta_2$ switch retains moments and the step counter, using the new value
in both the recurrence and native $1-\beta_2^{\mathrm{step}}$ correction.
Clipping acts on the complete accumulated raw gradient before Adam.
Caps target both MLP hidden affines, all 31 ResNet convolutions, and the two
QKV and four feed-forward Transformer affines. All 12 conditions per model
start at the same fork; no intervention waits for a later detected spike.

Probes retain the coverage in Table~\ref{tab:setup}.
The full training gradient is unchanged by probe coverage. Every state has loss
and motion measurements. FP32 training/storage and FP64 reductions are matched;
exact native replay checks qualify shared controls across workers.

We report the mean post-fork loss $\bar L$ and positive loss variation,
\begin{equation}
 \bar L=\frac{1}{e-f}\sum_{s=f+1}^{e}L_s,\qquad
 \mathrm{PV}=\sum_{s=f+1}^{e}\max(L_s-L_{s-1},0).
 \label{eq:vc-outcomes}
\end{equation}
PV accumulates upward loss changes, rather than counting excursions.
Final training loss, the last-32-state mean and full held-out losses retain
endpoint context. Held-out evaluations at $f$, $f+32$, midpoint, $e-32$ and $e$
are not used to select checkpoints or settings.

\subsubsection{Excursion suppression and training progress}
Figure~\ref{fig:vc-tradeoff} and Table~\ref{tab:vc-outcomes} separate loss
smoothness from learning progress. Reducing LR to $.1\eta$ removes
$99.6\%$, $99.0\%$, and $100\%$ of PV in MLP/ResNet/Transformer, but increases
ResNet mean loss by $89.5\%$. Late $\beta_2$ changes are model-dependent:
all three worsen MLP mean loss and PV, whereas $\beta_2=.9$ lowers
Transformer mean loss by $16.8\%$. Both clipping thresholds worsen mean loss
and PV in MLP and ResNet. Transformer clipping at $.5$ instead lowers them
by $17.8\%$ and $87.7\%$, respectively.

All nine cap/model combinations lower both mean loss and PV. The lowest observed
mean loss in each model belongs to a cap setting: q50 for MLP/Transformer
and q25 for ResNet. The preselected q50 reduces mean loss by
$32.2\%$, $9.4\%$, and $19.8\%$, with PV reductions of $87.8\%$, $2.8\%$, and $96.2\%$.
Quantile sensitivity is nonmonotone, and the finite-window mean does not establish
better terminal convergence. In particular, ResNet q50 increases final and
last-32 loss by $20.0\%$ and $21.5\%$; q25/q75 improve those endpoints.
The full trajectories and detailed outcomes are retained in the supplementary
comparison.

\subsubsection{Functional effect and limits}
Figure~\ref{fig:vc-motion} shows that q50 lowers the largest output RMS jump
by $83.2\%$, $23.1\%$, and $62.9\%$, while output reversals remain at
$89.2\%$, $66.9\%$, and $95.8\%$ of updates. Their baseline frequencies are
$99.3\%$, $36.9\%$, and $36.7\%$. Smaller extreme jumps therefore coexist with
persistent, sometimes more frequent, reversals. Total squared output motion
falls $88.6\%$ and $55.4\%$ in MLP and Transformer, but rises $39.7\%$
in ResNet. The effect concerns extreme displacements, not uniform suppression
of every measure of motion.

In the previously selected ResNet tracing window, steps $390$--$450$, q50
reduces the maximum loss from $.4407$ to $.1736$ ($60.6\%$).
Later oscillations explain why the complete $256$--$576$ window has only
$2.8\%$ lower PV. This secondary readout connects the intervention to the
traced event while retaining the full horizon as the primary comparison.

\begin{samepage}
For the squared-loss MLP, actual output motion has the exact accounting
\begin{equation}
 L_{s+1}-L_s=\frac{\langle F_s-Y,\Delta F_s\rangle}{N}
             +\frac{\|\Delta F_s\|_F^2}{2N},
 \label{eq:vc-output-loss}
\end{equation}
\end{samepage}
where $\Delta F_s=F_{s+1}-F_s$ and $N=5000$. For all 161 baseline and
70 q50 output-space loss rises above $10^{-8}$, the linear term is negative
and outweighed by the quadratic term. This is an exact output-space identity;
its linear term need not equal the parameter-space directional derivative.
Cross-entropy requires its own readout. The evidence supports displacement
control with persistent reversals, without identifying a unique causal source.

Training improvements also carry held-out tradeoffs: q50 terminal test loss
increases $0.53\%$, $2.42\%$, and $27.52\%$. These evaluations cover all 10,000
CIFAR-10 test images and 245,550 WikiText-2 test tokens; full validation and
five-state trajectories accompany the supplementary results. The implementation
has measurable cost: in the MLP q50 arm, controller time is $41.8$ ms/update,
compared with $4.4$ ms/full-gradient boundary. Instrumentation and host differences
preclude a time-to-loss claim. These single-initialization continuations use
unequal effective intervention strengths; reliability across initializations,
stochastic deployment and generalization benefits require additional evidence.
